\expandafter\def\csname HistoricalNoteText416\endcsname{An der Nadel einer Bussole, die sich an einer im Jahre 1709 verfertigten Sonnenuhr der hiesigen Sternwarte befindet, konnte 1841 durch neue Bestreichung bis zur Sättigung der Magnetismus nur auf das Dreifache erhöhet werden; an einer andern von 1603 nur auf das Fünffache. In der sehr wahrscheinlichen Voraussetzung, dass beide seit ihrer Verfertigung niemals neu gestrichen waren, und wenn man zugleich annimmt, dass sie ursprünglich auch bis zur Sättigung magnetisirt gewesen sind, und dass die Kraft allmählig in geometrischer Progression abgenommen hat, beträgt der jährliche Verlust bei der erstern $\frac{1}{120}$, bei der zweiten $\frac{1}{150}$, und noch weniger, falls die ursprüngliche Magnetisirung die Sättigung nicht erreicht hatte.}%
\SourceLeaf{412}zählten Winkel $A-P-p-B+Q$, und das aus der Einwirkung des Erdmagnetismus auf den Magnetstab entspringende von der Linken nach der Rechten positiv gerechnete Drehungsmoment ist $=RD\sin(A-P-p-B+Q)$. Wir haben mithin die Gleichung (1)
\[
0=\sin(A-P-p)+R\sin(A-P-p-B+Q)
\]
Wird der ganze Apparat aus der Gleichgewichtsstellung um den Winkel $z$ von der Rechten nach der Linken gedreht, so wirkt im entgegengesetzten Sinn das Drehungsmoment
\[
D\sin(z+A-P-p)+DR\sin(z+A-P-p-B+Q)
\]
welcher Ausdruck nach Entwicklung der beiden Sinus und unter Berücksichtigung der Gleichung (1) in
\[
D\sin z(\cos(A-P-p)+R\cos(A-P-p-B+Q))
\]
übergeht, also dem Sinus von $z$ proportional ist. Man hat also
\[
D(\cos(A-P-p)+R\cos(A-P-p-B+Q))
\]
wie die Directionskraft zu betrachten, die aus der Verbindung der statischen und magnetischen resultirt, und wir haben daher (2)
\[
\frac{S}{tt}=\cos(A-P-p)+R\cos(A-P-p-B+Q)
\]
Indem man in den beiden Gleichungen (1), (2) auf beiden Seiten quadrirt, und addirt, findet man (3)
\[
\frac{SS}{t^4}=1+2R\cos(Q-B)+RR
\]
Bezeichnet man mit $e$ die Basis der hyperbolischen Logarithmen und mit $i$ die imaginäre Einheit $\BareRoot{-1}$, so lassen sich die beiden Gleichungen (1), (2) bequem in Eine zusammenziehen
\[
\frac{S}{tt}=e^{i(A-P-p)}+Re^{i(A-P-p-B+Q)}
\]
oder noch einfacher in folgende (4)
\[
1=\frac{S}{tt}e^{i(P+p-A)}-Re^{i(Q-B)}
\]
\SourceLeaf{413}welche die beiden
\[
\begin{gathered}
1=\frac{S}{tt}\cos(P+p-A)-R\cos(Q-B)\\
\frac{S}{tt}\sin(P+p-A)=R\sin(Q-B)
\end{gathered}
\]
unter sich begreift.

Für die natürliche Lage, wo $Q=B$, wird
\[
tt=\frac{S}{1+R}
\]
für die verkehrte hingegen, wo $Q=B+180^0$,
\[
tt=\frac{S}{1-R}
\]

\Article{3}
Die transversale Stellung, im engern Sinne, erfordert, dass
\[
A-P-p-B+Q=\pm90^0
\]
wird, wo das obere Zeichen sich auf den Fall bezieht, wo der Nordpol des Magnetstabs auf der Westseite des magnetischen Meridians sein soll, das untere auf die östliche Lage. Es wird also nach (1)
\[
\sin(A-P-p)=\mp R
\]
Bezeichnet man demnach mit $\varphi$ den spitzen Winkel, dessen Sinus $=R$ ist, so wird für die westliche Stellung des Nordpols
\[
A=P+p-\varphi,\quad B=Q-\varphi-90^0
\]
für die östliche hingegen
\[
A=P+p+\varphi,\quad B=Q+\varphi+90^0
\]
Damit die Gleichgewichtsstellung dem durch den Lothfaden bezeichneten Skalenpunkte selbst entspreche, muss also die Spiegelalhidade auf $P-\varphi$ für den ersten Fall, und auf $P+\varphi$ für den zweiten gestellt werden.

Für die der Transversalstellung entsprechende Schwingungsdauer ergibt die Formel (2) sogleich
\SourceLeaf{414}\[
\frac{S}{tt}=\cos\varphi
\]
oder
\[
tt=\frac{S}{\BareRoot{(1-RR)}}
\]
Die Schwingungsdauer für die Transversalstellung ist demnach die mittlere Proportionale zwischen den Schwingungszeiten für die natürliche und für die verkehrte Stellung.

\Article{4}
Um klar übersehen zu können, in wiefern die Elemente veränderlich sind, müssen wir dieselben auf ihren Ursprung zurückführen.

Die Winkel $P$ und $Q$ sind jeder aus drei Theilen zusammengesetzt. Es besteht nemlich $P$ aus dem Winkel zwischen dem nach dem Nullpunkte des Kreises gehenden Radius und der die beiden untern Enden der Aufhängungsdrähte verbindenden geraden Linie; dem Winkel zwischen der die beiden obern Enden der Aufhängungsdrähte verbindenden geraden Linie und der optischen Axe des Fernrohrs (oder vielmehr zwischen den Projectionen dieser geraden Linie auf eine Horizontalebene, was auch bei allen andern Winkelschenkeln, die nicht selbst horizontal sind, oder unmittelbar einander nicht schneiden, stillschweigend verstanden wird); dem Winkel zwischen der Normale gegen den Spiegel und dem nach dem Hauptindex der Spiegelalhidade gehenden Radius.

Der erste Bestandtheil von $Q$ ist einerlei mit dem ersten Bestandtheile von $P$; der zweite ist der Winkel zwischen der die beiden obern Enden der Aufhängungsdrähte verbindenden geraden Linie und dem magnetischen Meridian; der dritte der Winkel zwischen der magnetischen Axe des im Schiffchen liegenden Magnetstabes und dem nach dem Hauptindex der Alhidade des Schiffchens gehenden Radius.

Alle diese fünf Winkel sind von der Linken nach der Rechten zu zählen. Es erhellt aus dieser Analyse, dass, insofern die Aufhängung des Instruments, die Verbindung des Spiegels mit seiner Alhidade und die Stellung des Beobachtungsfernrohrs unverrückt bleiben, $P$ ganz unveränderlich sein wird; dass aber $Q$ wegen seines zweiten Bestandtheils gerade dieselben Veränderungen erleidet, wie die magnetische Declination, die von der Linken nach der Rechten gehenden Veränderungen als positiv betrachtet.

\SourceLeaf{415}Die statische Directionskraft wird durch die Formel
\[
D=\frac{fgG}{4h}
\]
ausgedrückt, wo $G$ das Gewicht des Apparats (d. i. die durch die Schwerkraft multiplicirte Masse), $f$ den Abstand der Aufhängungsdrähte bei den untern, $g$ bei den obern Enden, $h$ die Höhe der obern Befestigung über der untern bedeutet; wenigstens insofern man die kleine Vergrösserung bei Seite setzt, welche jene Kraft noch durch die Reaction der einzelnen Aufhängungsdrähte gegen die Torsion erhält, was hier, wo zunächst nur von der Veränderlichkeit der ganzen Kraft die Rede ist, füglich geschehen kann. Bezeichnet man noch das Trägheitsmoment in Beziehung auf die verticale Drehungsaxe mit $K$, so wird, $\pi$ in üblicher Bedeutung genommen,
\[
S=\frac{\pi\pi K}{D}=\frac{4\pi\pi hK}{fgG}
\]
Es erhellt nun, dass die einzelnen Factoren $f,g,h,K$ in Folge des Temperaturwechsels Veränderungen erleiden, die freilich theils an sich sehr gering sind, theils wie weiter unten gezeigt werden wird, in dem Werthe von $S$ sich fast vollkommen compensiren. Als ganz unmerklich kann diejenige Ungleichheit angesehen werden, die aus dem ungleichen Gewichtsverlust in Folge ungleicher Luftdichtigkeit entspringt.

Die magnetische Directionskraft ist $=TM$, wenn $T$ die Intensität des horizontalen Erdmagnetismus, $M$ das Moment des Magnetismus im Magnetstabe ausdrückt; wir haben demnach
\[
R=\frac{TM}{D}=\frac{4TMh}{fgG}=\frac{STM}{\pi\pi K}
\]

Die Veränderlichkeit von $R$ beruht also auf einem dreifachen Grunde.

Erstlich auf der fortwährenden Veränderlichkeit von $T$; zweitens auf der Veränderlichkeit der Temperatur, welche nicht allein die Lineargrössen $f,g,h$ afficirt, sondern zugleich den Stabmagnetismus $M$; drittens auf der Veränderlichkeit von $M$ unabhängig von dem jedesmaligen Temperaturzustande.

In Beziehung auf die dritte Ursache sind unsere Kenntnisse bisher noch ziemlich unvollkommen. Bei den im 2. Bande der \textit{Resultate} mitgetheilten Versuchen des Hrn. Prof. \textsc{Weber} wurde der durch künstliche Erwärmung erlittene Verlust durch die nachherige Abkühlung niemals vollkommen ersetzt, sondern
\SourceLeaf{416}es blieb nach Wiederherstellung der anfänglichen Temperatur ein bedeutender nachhaltiger Verlust. Von der andern Seite lehrt die Erfahrung, dass Magnetnadeln ohne neue Bestreichung doch eine lange Reihe von Jahren, trotz der täglichen und jährlichen Abwechslung der Temperatur, einen bedeutenden Grad von Magnetismus behalten, woraus man also auf einen äusserst langsamen progressiven Verlust schliessen muss\HistoricalNote{416}. Es würde von grosser Wichtigkeit sein, die Bedingungen genau zu kennen, unter welchen der Temperaturwechsel den möglich kleinsten nachhaltigen Kraftverlust bewirkt. Ausser der Beschaffenheit und Härtung des Stahls, und einer kräftigen ursprünglichen Magnetisirung, wird es wahrscheinlich hauptsächlich darauf ankommen, dass seit dieser erst eine gewisse Zeit verflossen sein muss, dass die Temperaturänderungen gewisse Grenzen nicht überschreiten, und dass sie immer nur sehr langsam und allmählig erfolgen. Unter solchen Bedingungen wird es verstattet sein müssen, den magnetischen Zustand eines Magnetstabes --- wenn wir mit dieser Benennung sein auf eine bestimmte Normaltemperatur reducirtes magnetisches Moment bezeichnen --- während einer mässigen Zeit, z. B. einiger Tage, wie constant zu betrachten, und wenn nach einem längern Zeitraume eine entschiedene Abnahme gefunden wird, für die Zwischenzeit eine stetige Verminderung in geometrischer Progression zum Grunde zu legen. Die Ausführung des sinnreichen, von Hrn. Prof. \textsc{Weber} in dem weiter unten folgenden Aufsatze mitgetheilten Vorschlage scheint vorzüglich dazu geeignet, über diesen interessanten Gegenstand Licht zu verbreiten.

\Article{5}
Damit nun die Aufgabe, die Zahlenwerthe der Elemente eines Bifilarmagnetometers durch Versuche auszumitteln, eine präcise Bedeutung erhalte, verstehen wir unter den zu suchenden Werthen der veränderlichen Elemente diejenigen, die sich auf eine bestimmte Declination, eine bestimmte horizontale Intensität, eine bestimmte Temperatur und denjenigen magnetischen Zustand des
\SourceLeaf{417}Magnetstabes beziehen, welcher ihm zur Zeit dieser Versuche zukommt, wobei also die Veränderungen, welche letzterer nach längerer Zwischenzeit erleiden mag, gar nicht in Frage kommen. Wir bezeichnen diese Normalwerthe der veränderlichen Elemente mit $Q^0,R^0,S^0$ (indem $P$ schon für sich constant ist), und setzen allgemein
\[
Q=Q^0+q,\quad R=rR^0,\quad S=sS^0
\]
Auf gleiche Weise mögen $f^0,g^0,h^0,K^0,T^0,M^0$ die Normalwerthe der veränderlichen Grössen $f,g,h,K,T,M$ bezeichnen. Wir haben also sofort
\[
s=\frac{f^0g^0hK}{fgh^0K^0},\qquad r=\frac{f^0g^0hMT}{fgh^0M^0T^0}
\]

Um bei der Bestimmung der Elemente die während der dazu erforderlichen Operationen Statt findenden Veränderungen in der Richtung und Stärke der erdmagnetischen Kraft berücksichtigen zu können, muss natürlich ein Hülfsapparat zu Gebote stehen, am besten ein Unifilarmagnetometer, an welchem gleichzeitig Schwingungsdauer und Stand beobachtet werden. Zugleich dient dieses Hülfsmagnetometer dazu, die zu wählende Normaldeclination und Normalintensität nachweisbar zu machen, zunächst dadurch, dass man jene einem bestimmten Skalenpunkte, diese einer bestimmten Schwingungsdauer für die Normaltemperatur entsprechen lässt, wobei man dann auch in seiner Gewalt hat, beide Normalgrössen nach bekannten Methoden auf absolutes Maass zu bringen. Hiernach ist ohne weiteres $q$ der in Bogentheile verwandelte Unterschied des am Hülfsmagnetometer beobachteten Standes vom Normalstande. Bezeichnet man ferner, was am Bifilarmagnetometer $M,K,t$ ist, für das Hülfsmagnetometer mit $m,k,\theta$, und die Normalwerthe dieser Grössen mit $m^0,k^0,\theta^0$, so wird
\[
\frac{\theta\theta}{\theta^0\theta^0}=\frac{km^0T^0}{k^0mT}
\]
und folglich
\[
r=\frac{f^0g^0hkm^0M\theta^0\theta^0}{fgh^0k^0mM^0\theta\theta}=\frac{skK^0m^0M\theta^0\theta^0}{k^0KmM^0\theta\theta}
\]

Von den sieben Factoren $\frac{f^0}{f},\frac{g^0}{g},\frac{h}{h^0},\frac{k}{k^0},\frac{K^0}{K},\frac{m^0}{m},\frac{M}{M^0}$, welche in den Ausdrücken für $s$ und $r$ vorkommen, wird man die fünf ersten nach der Ausdehnung, welche die betreffenden Metalle durch die Temperatur erleiden, die beiden letzten hingegen nach der besten Kenntniss, die man vom Einfluss der Tem\SourceLeaf{418}peratur auf den Stabmagnetismus besitzt, zu berechnen haben, indem das, was wir den magnetischen Zustand genannt haben, bei beiden Magnetstäben während der hier in Rede stehenden Operationen wie constant betrachtet wird. Wir fügen in Beziehung auf diese Rechnung noch einige Entwickelungen bei.

Indem wir zur Normaltemperatur den Gefrierpunkt wählen, bezeichnen wir mit $c$ und $c'$ die Temperatur im Kasten des Bifilar- und des Hülfsmagnetometers, mit $c''$ die Temperatur bei der obern Befestigung der Aufhängungsdrähte des erstern; ferner, für Einen Grad Wärmezunahme, die Ausdehnung des Stahls mit $\alpha$, des Messings mit $\beta$, und die Abnahme des Stabmagnetismus für die Stäbe der beiden Apparate mit $\gamma$ und $\gamma'$. Da die Veränderung des Trägheitsmoments der beiden Apparate dem bei weitem grössten Theile nach von der Ausdehnung der Magnetstäbe selbst herrührt, so wird man ohne Bedenken
\[
\frac{K}{K^0}=(1+\alpha c)^2,\qquad \frac{k}{k^0}=(1+\alpha c')^2
\]
setzen; für die Ausdehnung der Aufhängungsdrähte, wenn sie, wie am hiesigen Apparate, Stahldrähte sind, wird man denselben Coëfficienten $\alpha$ beibehalten, und für ihre Temperatur $\tfrac{1}{2}(c+c'')$ annehmen können, so dass
\[
\frac{h}{h^0}=1+\tfrac{1}{2}\alpha(c+c'')
\]
wird. Wir haben mithin (1)
\[
s=\frac{(1+\alpha c)^2(1+\tfrac{1}{2}\alpha(c+c''))}{(1+\beta c)(1+\beta c'')}
\]
wofür man auch, hinlänglich genau,
\[
s=1+(3\alpha-2\beta)c-(\beta-\tfrac{1}{2}\alpha)(c''-c)
\]
setzen kann. Da nun, der Erfahrung zufolge, sehr nahe $\beta=\tfrac{3}{2}\alpha$ ist, so wird, sehr nahe, (2)
\[
s=1-\alpha(c''-c)
\]
d i. die Veränderung des Elements $S$ ist nur von der Ungleichheit der untern und obern Temperatur abhängig, so dass in der Regel $S$ wie ganz constant betrachtet werden kann.

Wir haben ferner (3)
\[
r=\frac{s\theta^0\theta^0}{\theta\theta}.\frac{1-\gamma c}{1-\gamma'c'}.\left(\frac{1+\alpha c'}{1+\alpha c}\right)^2
\]
\SourceLeaf{419}oder wenn die Temperaturänderungen auf beide Stäbe gleichen Einfluss haben, d. i. wenn $\gamma'=\gamma$ ist, hinlänglich genau,
\[
r=\frac{s\theta^0\theta^0}{\theta\theta}(1+(\gamma+2\alpha)(c'-c))
\]
oder in Gemässheit von (2), eben so genau (4)
\[
r=\frac{\theta^0\theta^0}{\theta\theta}(1+(\gamma+2\alpha)(c'-c)-\alpha(c''-c))
\]

Endlich muss noch der Umstand bemerkt werden, dass durch die Vergleichungsbeobachtungen am Unifilarmagnetometer nicht der für einen bestimmten Augenblick geltende Werth von $\frac{T}{T^0}$ abgeleitet werden kann, sondern nur der Mittelwerth für die ganze Zeit, welche die Schwingungsbeobachtungen umfassen. Es versteht sich also von selbst, dass auch alle die andern Grössen, mit denen jene Schwingungsbeobachtungen als gleichzeitige unmittelbar oder mittelbar combinirt werden sollen, sich gleichfalls als Mittelwerthe auf denselben Zeitraum beziehen müssen.

\Article{6}
Die kunstloseste Art, die vier Elemente auszumitteln, ist folgende:

Bei willkürlicher Stellung des Schiffchens legt man anstatt des Magnetstabes einen nicht magnetischen Stab, ungefähr von gleichem Gewicht, in dasselbe, und gibt dem Spiegel eine solche Stellung, dass in der Gleichgewichtslage das Bild irgend eines Punkts der Skale auf dem Fadenkreuz des Beobachtungsfernrohrs erscheint, wo dann, $A$ und $p$ in der obigen Bedeutung genommen, $P=A-p$ wird. Um das Resultat von einer sehr genauen Kenntniss des Werthes der Skalentheile oder von einer sehr scharfen Reduction derselben auf Bogentheile unabhängiger zu machen, mag man die Operation, wenn das erstemal $p$ noch sehr gross ausgefallen ist, mit einer neuen sehr genäherten Stellung des Spiegels wiederholen. Am meisten geeignet für diese Operation ist ein mit Blei belasteter Holzstab; das ungefähr gleiche Gewicht wird deswegen erfordert, um eine kleine Torsion, welche bei der Gleichgewichtsstellung des Ganzen die Aufhängungsdrähte für sich genommen möglicherweise haben könnten, unschädlich zu machen.

Ohne nun die Stellung des Spiegels weiter zu ändern, legt man anstatt der vorigen Belastung den Magnetstab in das Schiffchen, welches dann so gestellt
\SourceLeaf{420}werden soll, dass dem Ruhestande derselbe Skalenpunkt entspreche, wie zuletzt bei der nicht magnetischen Belastung. Man gelangt dazu, indem man durch Versuche zwei verschiedene Stellungen des Schiffchens ermittelt, zwischen welche die gesuchte fällt, und auf die bei jenen sich ergebenden Ablesungen an der Skale ein einfaches Interpolationsverfahren anwendet. Man kann sich hiebei entweder der natürlichen oder der verkehrten Lage des Magnetstabes bedienen; im ersten Falle ist das sich für $B$ (die Stellung der Alhidade des Schiffchens) ergebende Resultat $=Q$, im zweiten $=Q\pm180^0$. Die Anwendung der verkehrten Lage hat den Vorzug grösserer Schärfe, weil einer kleinen Änderung von $B$ eine grosse Änderung der Skalentheile entspricht, die Anwendung der natürlichen Lage hingegen ist in so fern etwas bequemer, als man dabei dem Schiffchen eine nicht über die Grenzen der Skale hinausgehende Lage leichter geben kann. Man thut daher wohl, zur Vermeidung beschwerlichen Herumtastens, mit der natürlichen Lage anzufangen, das gefundene Resultat aber nur wie eine Vorbereitung zu betrachten, um bei den Versuchen in verkehrter Lage auf zwei nahe zusammenliegende Theilstriche einstellen zu können.

Das gefundene Resultat für $Q$ bezieht sich auf diejenige Lage des magnetischen Meridians, welche derselbe in oder zwischen den beiden letzten Versuchen gehabt hat, und mehr als eine solche schwankende Bestimmung ist nicht zu fordern, wenn man keinen Hülfsapparat zu vergleichenden Beobachtungen anwenden kann. Steht aber ein Hülfsapparat zu Gebote, so geben gleichzeitige Standbeobachtungen an demselben die jenen beiden Beobachtungen correspondirenden Werthe von $q$ und das obige Interpolationsverfahren auf die beiden Werthe von $B-q$ angewandt ergibt dann den Werth von $Q^0$ oder $Q^0\pm180^0$.

Endlich beobachtet man die Schwingungsdauer sowohl in der natürlichen als in der verkehrten Lage; man stellt zu dem Ende die Alhidade des Schiffchens so genau man kann auf denjenigen Werth von $Q$ (und beziehungsweise von $Q+180^0$), der eben beim Anfang der Schwingungsbeobachtungen gilt. Die Schwingungsdauer in der natürlichen Lage sei $t'$, in der verkehrten $t''$; kann man gleichzeitig Schwingungen am Hülfsmagnetometer beobachten, so erhält man dadurch die correspondirenden Werthe von $r$, die mit $r',r''$ bezeichnet werden mögen; will man auch die Veränderlichkeit von $S$ berücksichtigen, so mögen $s',s''$ die correspondirenden Werthe von $s$ sein. Die kleinen Veränderungen in der Lage des magnetischen Meridians während der Schwingungsbeobachtungen
\SourceLeaf{421}werden in der Regel keinen merklichen Einfluss auf die Resultate haben. Die beiden Gleichungen am Schluss des 2. Artikels werden demnach
\[
\begin{aligned}
t't'&=\frac{s'S^0}{1+r'R^0}\\
t''t''&=\frac{s''S^0}{1-r''R^0}
\end{aligned}
\]
woraus durch Elimination folgt
\[
\begin{aligned}
R^0&=\frac{s't''t''-s''t't'}{r''s't''t''+r's''t't'}\\
S^0&=\frac{(r'+r'')t't't''t''}{r''s't''t''+r's''t't'}
\end{aligned}
\]

Nach der im 5. Art. gemachten Bemerkung kann man füglich $S$ wie constant betrachten, oder $s'=s''=1$ setzen, wodurch die Formeln in
\[
\begin{aligned}
R^0&=\frac{t''t''-t't'}{r''t''t''+r't't'}\\
S&=\frac{(r'+r'')t't't''t''}{r''t''t''+r't't'}
\end{aligned}
\]
übergehen. Kann man aber keine Vergleichungsbeobachtungen an einem Hülfsapparat zuziehen, so bleibt nichts übrig, als geradezu
\[
\begin{aligned}
R&=\frac{t''t''-t't'}{t''t''+t't'}\\
S&=\frac{2t't't''t''}{t''t''+t't'}
\end{aligned}
\]
zu setzen, und es ist klar, dass der so gefundene Werth von $R$ nur eine Art von Mittel zwischen den für die beiden Schwingungssätze geltenden bedeuten, $S$ aber mit einer kleinen von der Ungleichheit der letztern abhängenden Unrichtigkeit behaftet bleiben wird.

\Article{7}
Die allgemeinere Auflösung unsrer Aufgabe gründen wir auf die gleichzeitigen Beobachtungen von Schwingungsdauer und Gleichgewichtsstand des Bifilarmagnetometers bei zwei beliebigen ungleichen Stellungen des Schiffchens. Wir bezeichnen die bestimmten Werthe der Grössen $A,B,p,Q,R,S,t$

für den ersten Satz der Beobachtungen mit $A',B',p',Q^0+q',r'R^0,s'S^0,t'$;

für den zweiten Satz mit $A'',B'',p'',Q^0+q'',r''R^0,s''S^0,t''$.

